\begin{problem}{Подсчёт ворон}{стандартный ввод}{стандартный вывод}{1 секунда}{256 мегабайт}

Вася изучает ворон и обратил внимание, что они сидят на деревьях, расположенных вдоль улицы, номера которых кратны числу $C$ (т.е. на деревьях с номерами $C$, $2C$, $3C$, ...). Вася шел мимо деревьев с номерами от $A$ до $B$ включительно и считал ворон. Сколько ворон он насчитал?

\InputFile
Во входных данных записаны три натуральных числа $A$, $B$ и $C$ ($1 \le A \le B \le 2 \cdot 10^9$, $1 \le C \le 2 \cdot 10^9$).

\OutputFile
Выведите одно число --- количество ворон.

\Example

\begin{example}
\exmpfile{example.01}{example.01.a}%
\end{example}

\Note
В примере вороны сидят на деревьях $1850$ и $1900$.

\end{problem}

